\documentclass[12pt]{article}
\usepackage{amsmath, amssymb}
\usepackage{geometry}
\usepackage{hyperref}
\geometry{margin=1in}

\title{On the Nature of Mass:\\
Complexity Generation and Spacetime Curvature}
\author{Paul A. I. Reilly}
\date{[DRAFT v0.9 --- not for circulation]}

\begin{document}

\maketitle

%% ============================================================
\begin{abstract}
%% ============================================================

We show that the radial component of the exterior Schwarzschild metric
can be derived exactly from one central hypothesis and four supporting
assumptions, without invoking Einstein's equations.
The central hypothesis is that mass is a rate of irreversible information
generation: an isolated mass $M$ of pressureless dust produces classical
information at rate $dC/d\tau = \alpha Mc^2/\hbar$, where $\alpha$ is a
single free parameter.
Extensions to matter with pressure, shear, or momentum flux are addressed
in the companion paper (Paper~II).
Four assumptions govern how this information is geometrically accommodated:
(1) information propagates at the speed of light; (2) the vacuum near
matter operates at holographic saturation; (3) each nat of information
requires a minimum of four Planck areas (Bekenstein); and (4) nature is
parsimonious.
From these, the radial metric component $g_{rr}$ follows algebraically.
Matching to the known Schwarzschild solution fixes $\alpha = 2\pi$; this
value is independently confirmed by the Bekenstein bound at the Compton
scale, establishing the hypothesis as a consistency requirement rather
than an ansatz.
We further present a well-motivated argument for the temporal component
$g_{tt}$, identifying one step as a conjecture open to community
scrutiny; an independent derivation of $g_{tt}$ appears in the companion
paper via the Einstein equations, themselves derived from Screen Theory
premises.
The argument applies equally to isolated masses of any scale.
A cosmological motivation for the holographic saturation assumption,
based on an Olbers-style integration of information flux from all matter
in the observable universe, is developed in Appendix~\ref{app:A2}.

\end{abstract}

%% ============================================================
\section{Introduction}
%% ============================================================

\subsection{The Central Claim}

Mass is ordinarily defined operationally: it measures resistance to
acceleration and determines the strength of gravitational interaction.
In this paper we propose a more fundamental characterization.

\emph{Mass is the rate at which a system generates irreversible classical
information.}

Spacetime curvature is the geometric accommodation required by this
ongoing process. More precisely, we propose the hypothesis
\begin{equation}
  \frac{dC}{d\tau} = \alpha \frac{Mc^2}{\hbar},
  \label{eq:hypothesis}
\end{equation}
where $C$ denotes the amount of irreducible classical information
generated by an isolated mass $M$, and $\alpha$ is a dimensionless
constant determined by consistency conditions.
In the derivation below we show that the radial component of the exterior
Schwarzschild metric follows directly from this hypothesis together with
four physical assumptions governing how information is geometrically
accommodated.
Hypothesis~\eqref{eq:hypothesis} is stated for pressureless dust;
extensions to matter with pressure, shear, or momentum flux are more
complex and are addressed in the companion paper (Paper~II).

In this framework classical information is not conserved but increases
monotonically through irreversible physical processes.
The accumulation of such information defines a natural informational
arrow of time.
Spacetime geometry responds not merely to energy but to the irreversible
narrowing of the set of physically possible futures.

The quantity $C$ refers to the irreducible description length of the
sequence of irreversible events produced by physical processes.
Operationally this corresponds to the Kolmogorov complexity
\cite{Kolmogorov1965} of the history of objective state reductions,
evaluated with respect to a description language containing the known
laws of physics.
Unitary quantum evolution does not increase this complexity because the
future state is algorithmically determined by the past.
Only irreversible transitions introduce genuinely new constraints on the
future evolution of the universe.

The argument originates in a simple observation: $E = hf$ and $E = Mc^2$
share a left-hand side.
Setting them equal gives every isolated mass a characteristic frequency
$f = Mc^2/h$ --- the Compton frequency --- and raises a concrete
question: what is the mass doing at that frequency?
The answer proposed here is that it behaves as a Shannon source
\cite{Shannon1948}, producing irreversible classical information at that
rate.
The holographic principle then supplies the geometric link: information
requires physical expression as boundary area, and since it propagates
outward at $c$, area accumulates in the surrounding region.
The key physical insight is that information generated near the mass
spends time proportional to $r/c$ \emph{in transit} before reaching
a bounding surface at radius $r$.
More area accumulates per unit radius closer to the source, giving
gravity its $1/r$ character and producing the Schwarzschild radial
metric.

A note on terminology: throughout this paper we use \emph{isolated mass}
rather than \emph{particle}.
The argument applies equally to a proton, a neutron star, or the
observable universe.
The Compton scale $\lambda_C = \hbar/Mc$ is the relevant length scale
in each case, regardless of internal structure.

\subsection{Irreversible Processes and Information Generation}
\label{sec:irreversible}

Quantum dynamics contains two qualitatively different forms of
evolution \cite{Penrose1994}: reversible unitary evolution (U-process),
which preserves information content, and irreversible state reduction
(R-process), which produces definite classical outcomes and thereby
generates new classical information.

A simple example illustrates the principle.
Consider a 50/50 beam splitter illuminated by single photons at a rate
of $10^9$ per second.
Each photon is detected in either the transmitted or reflected arm,
producing a binary sequence of outcomes.
Because the events are statistically independent, this sequence contains
roughly $10^9$ bits of algorithmically incompressible classical
information after one second.
The mass equivalent implied by hypothesis~\eqref{eq:hypothesis} is
\begin{equation}
  M = \frac{\hbar}{2\pi c^2}\frac{dC}{d\tau}
    = \frac{1.055\times10^{-34}}{2\pi \times 9\times10^{16}}
      \times 10^9
    \approx 1.86\times10^{-43}\;\text{kg},
\end{equation}
confirming that the associated geometric effect is undetectably small
under ordinary conditions.
The example demonstrates that irreversible quantum processes generate
objective classical information at well-defined rates.
The framework asserts that matter continuously drives analogous processes
through its interaction with vacuum degrees of freedom at a rate set by
the Compton frequency.
The microscopic mechanism is not specified here and remains an open
question.

The microscopic mechanism responsible for this information generation
is not specified here. It will be explored in companion papers,
where we develop a concrete model of the matter-vacuum interaction
that produces the proposed irreversible information at the Compton rate.
The present paper remains agnostic about this mechanism, focusing
instead on the kinematic consequences for spacetime geometry.

A potential confusion deserves direct address.
The complexity generation rate $dC/d\tau$ for one kilogram of matter is
vastly larger than the thermodynamic entropy of the same mass
($\sim 10^{25}$ nats for a kilogram of gas at standard conditions).
These are different ledgers.
Thermodynamic entropy counts the arrangement of matter's own microscopic
degrees of freedom.
The complexity generation rate counts the ongoing output of the
matter-vacuum interaction --- information flowing into the structure of
the surrounding vacuum, not into the internal state of the matter.
The large ratio is a consistency requirement rather than a liability:
the derivation produces the correct Schwarzschild geometry precisely
because the generation rate is set by mass-energy (via $\hbar$ and
$\alpha = 2\pi$), not by the thermodynamic state of the matter.
A rate comparable to thermodynamic entropy would produce no measurable
curvature.

\subsection{Relation to Prior Work}
\label{sec:prior}

The identification of quantum state reduction as the source of
irreversibility follows Penrose \cite{Penrose1994, HawkingPenrose1996},
who distinguished unitary from irreversible quantum evolution and
proposed that R-processes are connected to spacetime geometry.
Our framework inverts Penrose's causal direction: where Penrose proposed
that geometry \emph{drives} state reduction, we find that irreversible
transitions \emph{produce} geometry.
The R-process is primary; curvature is its accommodation.
We remain agnostic about the physical mechanism responsible for the
transition itself.

The holographic principle --- that the information content of a region
is bounded by its boundary area --- provides the bridge from information
generation to geometry.
Developed by 't~Hooft \cite{tHooft1993} and Susskind \cite{Susskind1995}
and given precise form by Bekenstein's bound \cite{Bekenstein1973,
Bekenstein1981}, it is the second pillar of our argument.
That entanglement entropy in quantum field theory scales with boundary
area rather than volume \cite{Srednicki1993} provides concrete evidence
that area-based information structure is a natural feature of quantum
spacetime, independently supporting assumption A2 below.

Several previous approaches have explored connections between gravity
and information.
The present framework differs from thermodynamic derivations of
gravitational dynamics in a key respect: rather than treating gravity
as an emergent entropic phenomenon, we identify mass itself as a rate
of irreversible information generation and interpret curvature as the
kinematic accommodation of that production.
The companion papers develop this distinction in detail.

\subsection{Hypothesis and Assumptions}
\label{sec:hypothesis}

The derivation rests on one central hypothesis and four assumptions.

\medskip
\noindent\textbf{H (Generation hypothesis).}
An isolated mass $M$ of pressureless dust generates irreversible
classical information at rate
$dC/d\tau = \alpha Mc^2/\hbar$,
where $\alpha$ is a dimensionless constant determined by internal
consistency with the assumptions below.

\medskip
\noindent\textbf{A1 (Propagation).}
Information propagates at the speed of light.

\medskip
\noindent\textbf{A2 (Saturation).}
The vacuum near matter operates at holographic saturation: newly
generated information requires a corresponding increase in bounding
area.
A2 is an assumption of this framework; Srednicki \cite{Srednicki1993}
provides independent theoretical support from the area scaling of
entanglement entropy in quantum field theory, and Appendix~\ref{app:A2}
provides an independent quantitative motivation from an Olbers-style
integration of cosmological information flux.
The successful derivation of the Schwarzschild metric in Section~2
provides a posteriori consistency, though not proof.

\medskip
\noindent\textbf{A3 (Bekenstein bound).}
Each nat of information requires a minimum area of $4\ell_P^2$, where
$\ell_P = \sqrt{G\hbar/c^3}$.

\medskip
\noindent\textbf{A4 (Minimality).}
Nature produces no excess area beyond what the information requires.

\subsection{Departures from Standard Frameworks}
\label{sec:departures}

The central hypothesis has several consequences that depart from
standard physical assumptions.
We state them plainly here so that the reader can assess them
explicitly rather than encounter them as surprises.

\medskip
\noindent\textbf{Quantum evolution is not strictly unitary.}
R-processes are genuinely irreversible, not merely apparent
irreversibility arising from coarse-graining or loss of phase
information.
The framework treats state reduction as a primitive, not as something
to be derived from or reconciled with unitary evolution.
This is a deliberate departure from the Everettian programme; it
follows Penrose's insistence that R-processes are physically real,
though with the causal direction reversed.

\medskip
\noindent\textbf{Information is not conserved in the sense of a
fixed total.}
Each R-process produces a classical, objective outcome: a settled
fact added permanently to the universe's record.
In the sense of a museum conservator, information is never lost ---
settled outcomes are not undone --- but the collection is always
growing.
The total classical information content of the universe increases
monotonically.
This monotonic accumulation is the physical origin of the arrow of
time in the present framework.
Irreversibility deepens in stages: a measurement outcome is
irreversible in the practical sense (phase information gone);
information scrambled with many degrees of freedom is irreversible
in the thermodynamic sense; but the deepest form is causal ---
information crossing a horizon is placed permanently beyond the reach
of every future observer within the accessible universe.
Horizon-crossing is the universe's irrevocable commitment.

\medskip
\noindent\textbf{CPT symmetry is broken at the level of fundamental
dynamics.}
If R-processes are genuinely time-asymmetric, CPT violation is not
merely a thermodynamic artifact but a feature of the underlying
physics.
The expected magnitude, suppressed by approximately
$1/(\text{age of universe}) \approx 10^{-17}$--$10^{-18}$~s$^{-1}$,
corresponds to an energy scale of order $10^{-42}$~GeV --- roughly
24 orders of magnitude below the current best experimental bounds
from neutral kaon systems ($\sim 10^{-18}$~GeV \cite{PDG2024}).
The prediction is therefore consistent with all current experimental
results and not accessible to near-term experiments.
Whether the suppression is uniform or concentrates near compact masses
where the local generation rate is high is a question for companion
papers and future experiment.

\medskip
\noindent\textbf{Spacetime curvature is kinematic rather than
dynamical.}
In standard general relativity, curvature is governed by field
equations: it is the dynamical response of spacetime to an
energy-momentum source.
In the present framework, curvature is a kinematic consequence of
information-theoretic constraints --- the geometric shape that
information generation necessarily imposes on the surrounding
spacetime, in the same way that Lorentz contraction is a kinematic
consequence of the invariance of $c$ rather than a dynamical force.
Whether this reframing has observable consequences beyond reproducing
the standard results is a question the companion papers begin to
address.

\subsection{Scope of the Present Paper}
\label{sec:scope}

The aim of the present paper is limited.
We show that the radial component of the exterior Schwarzschild metric
follows directly from H and A1--A4 without invoking Einstein's field
equations.
An argument for the temporal component is presented as a conjecture
open to further scrutiny.
Questions concerning the microscopic origin of the information
generation process, the derivation of the Einstein equations, and
cosmological information balance are addressed in companion papers.

%% ============================================================
\section{Derivation of the Exterior Schwarzschild Metric}
%% ============================================================

We derive the exterior Schwarzschild metric from hypothesis H and
assumptions A1--A4, without invoking Einstein's field equations.

\subsection{Area Response to Information Growth}
\label{sec:area-response}

The derivation rests on the following chain of reasoning.
Classical information is generated through irreversible quantum state
reductions.
Because such reductions are non-unitary, classical information
increases monotonically, providing a natural informational arrow of
time.
Information contained within a spatial region is subject to the
Bekenstein--Hawking bound: the maximum entropy associated with a
region is proportional to the area of its bounding surface,
\begin{equation}
  S_{\max} = \frac{A}{4\ell_P^2}.
  \label{eq:bekenstein-hawking}
\end{equation}
If classical information within a region increases through irreversible
state reductions, the holographic bound implies that the bounding area
must increase accordingly.
This requirement links information growth directly to geometry.

A preliminary calculation shows that the ratio of holographic
information to holographic capacity yields the exact Schwarzschild
parameter.
The Bekenstein bound for mass $M$ within radius $r$ (in natural units
$\hbar = c = 1$) is
\begin{equation*}
  S_{\rm Bekenstein} = 2\pi Mr.
\end{equation*}
The holographic capacity of the bounding sphere of radius $r$
(with $\ell_P^2 = G$ in natural units) is
\begin{equation*}
  S_{\max} = \frac{4\pi r^2}{4G} = \frac{\pi r^2}{G}.
\end{equation*}
The fractional saturation is therefore, exactly,
\begin{equation*}
  \frac{S_{\rm Bekenstein}}{S_{\max}}
  = \frac{2\pi Mr}{\pi r^2/G}
  = \frac{2GM}{r}
  = \frac{r_S}{r},
\end{equation*}
where $r_S = 2GM/c^2$ is the Schwarzschild radius.
This is not a scaling result: the numerical prefactors work out
exactly, with no residual dimensionless factor.
The gravitational field strength at radius $r$ is precisely the
degree to which the region's informational content saturates its
holographic capacity.
It is exactly this ratio $r_S/r$ that controls every deviation from
flat space in both components of the Schwarzschild metric:
$g_{tt} = -(1 - r_S/r)$ and $g_{rr} = (1 - r_S/r)^{-1}$.

For a general static spherically symmetric line element
\cite{Wald1984}
\begin{equation}
  ds^2 = -A(r)\,dt^2 + B(r)\,dr^2 + r^2\,d\Omega^2,
  \label{eq:general-metric}
\end{equation}
the radial metric coefficient $B(r)$ controls the physical length
associated with a coordinate radial displacement.
The derivation in Sections~2.3 and~2.4 makes the informational
origin of $B(r)$ precise.

\subsection{Setup and Coordinate Conventions}
\label{sec:setup}

Consider a spherically symmetric, stationary mass $M$ of pressureless
dust.
We work in the exterior vacuum region $r > R$, where $R$ is the radius
of the matter distribution, and the geometry approaches flat Minkowski
spacetime as $r \to \infty$.
We adopt Schwarzschild coordinates $(t, r, \theta, \phi)$ \cite{Wald1984}.
The Schwarzschild radius is
\begin{equation}
  r_S = \frac{2GM}{c^2}.
  \label{eq:rs}
\end{equation}
This is a geometric length scale; it does not imply the object is a
black hole.
We proceed in two steps: $g_{rr}$ from area excess (Section~2.3),
then the argument for $g_{tt}$ (Section~2.4).

\subsection{The Radial Component: Area Excess from Information
in Transit}
\label{sec:grr}

By H, the central mass generates information at rate
\begin{equation}
  \frac{dC}{d\tau} = \alpha \frac{Mc^2}{\hbar}.
  \label{eq:gen-rate}
\end{equation}
By A1, information propagates outward at $c$.
In the stationary configuration, the information currently in transit
between the source and a sphere at coordinate radius $r$ was generated
during the preceding proper time interval $r/c$, giving
\begin{equation}
  \Delta C = \frac{dC}{d\tau}\cdot\frac{r}{c}
           = \frac{\alpha Mcr}{\hbar}.
  \label{eq:transit}
\end{equation}
By A2 (saturation), this information requires geometric expression as
area.
By A3 (Bekenstein), each nat requires $4\ell_P^2$ of area.
By A4 (minimality), the area excess is exactly
\begin{equation}
  \Delta A = 4\ell_P^2\,\Delta C
           = 4\frac{G\hbar}{c^3}\cdot\frac{\alpha Mcr}{\hbar}
           = \frac{4\alpha GMr}{c^2},
  \label{eq:deltaA}
\end{equation}
which simplifies to
\begin{equation}
  \Delta A = 2\alpha\, r_S r.
  \label{eq:deltaA-simple}
\end{equation}
We define the radial coordinate $r$ as the areal coordinate: the
physical area of a sphere at coordinate $r$ is $A_{\rm phys} = 4\pi r^2$.
The area that sphere would have absent the mass's information
contribution is
\begin{equation}
  A_{\rm unstretched} = 4\pi r^2 - 2\alpha r_S r.
  \label{eq:unstretched}
\end{equation}
The radial metric component is the ratio of physical to unstretched
area:
\begin{equation}
  g_{rr} = \frac{4\pi r^2}{4\pi r^2 - 2\alpha r_S r}
          = \frac{1}{1 - \alpha r_S/2\pi r}.
  \label{eq:grr-alpha}
\end{equation}
The area contribution from information in transit grows linearly with
$r$ --- exactly the profile required for a $1/r$ gravitational potential.
Matching to the known Schwarzschild metric $g_{rr} = (1 - r_S/r)^{-1}$
fixes $\alpha = 2\pi$, so
\begin{equation}
  g_{rr} = \frac{1}{1 - r_S/r}.
  \label{eq:grr}
\end{equation}
This is the radial component of the Schwarzschild metric, derived from
information accounting alone.
Each step follows from H and A1--A4 in sequence.
No field equations, curvature tensors, or prior assumptions about
geometry are required.

\subsection{The Temporal Component: A Conjecture and an Invitation}
\label{sec:gtt}

The derivation of $g_{rr}$ is kinematic: it follows from counting
information in transit.
The temporal component $g_{tt}$ requires a further argument.
We present the strongest argument we have, identify the step that
remains a conjecture, and invite community scrutiny.
An independent derivation of $g_{tt}$ via the Einstein equations
appears in Paper~II.

The configuration is stationary, so the spacetime admits a timelike
Killing vector $\partial/\partial t$.
The energy of any null ray --- the carrier of information --- is
conserved along each trajectory in the exterior:
\begin{equation}
  E = -g_{tt}\frac{dt}{d\lambda} = \text{const along each null ray.}
  \label{eq:killing-energy}
\end{equation}
In the source-free exterior, no information is generated.
The integrated information flux through any concentric sphere is
therefore constant:
\begin{equation}
  \Phi(r) = \text{const} \qquad \text{for all } r > R.
  \label{eq:flux-conservation}
\end{equation}
The most general static spherically symmetric metric takes the form
\cite{Wald1984}
\begin{equation}
  ds^2 = -A(r)\,dt^2 + B(r)\,dr^2 + r^2\,d\Omega^2,
  \label{eq:metric-general}
\end{equation}
with $B(r) = (1 - r_S/r)^{-1}$ already determined by Section~2.3.
We write
\begin{equation}
  A(r) = e^{2\Phi(r)}\left(1 - \frac{r_S}{r}\right)
  \label{eq:Ar}
\end{equation}
for a free function $\Phi(r)$.
Asymptotic flatness requires $\Phi(r) \to 0$ as $r \to \infty$;
the general expansion is $\Phi(r) = \beta/r + \gamma/r^2 + \cdots$
Each nonzero coefficient introduces a free parameter that must be
sourced by some physical field.
In the present setup --- vacuum exterior, single length scale $r_S$
already fixed, no additional fields --- no such source is present.

\medskip
\noindent\textbf{Conjecture:} $\Phi(r) = 0$ identically.

\medskip
By A4 (minimality), no geometric structure is produced beyond what
the information constraints require.
Any nonzero $\Phi(r)$ introduces structure sourced by nothing in the
physical setup, violating A4.
Therefore $g_{tt} = -(1 - r_S/r)$.

\medskip
\noindent\textbf{Where the argument needs scrutiny.}
A4 is doing substantial work here.
In GR, $\Phi = 0$ follows from the vacuum Einstein equations via
Birkhoff's theorem --- a rigorous result.
We reach the same conclusion by a different route.
The open problem is to show rigorously that A4 excludes all unsourced
free parameters in the temporal component, without implicitly assuming
the conclusion.
We present this as a conjecture; the companion paper (Paper~II) provides
an independent derivation that may serve as confirmation pending a
direct proof.

\subsection{The Exterior Schwarzschild Metric}
\label{sec:metric}

If the conjecture of Section~2.4 holds, the complete line element in
the exterior region $r > R$ is
\begin{equation}
  ds^2 = -\!\left(1 - \frac{r_S}{r}\right)dt^2
         + \left(1 - \frac{r_S}{r}\right)^{-1}dr^2
         + r^2\!\left(d\theta^2 + \sin^2\!\theta\,d\phi^2\right),
  \label{eq:schwarzschild}
\end{equation}
with $r_S = 2GM/c^2$ and $\alpha = 2\pi$.
The radial component is exact.
The temporal component depends on the conjecture of Section~2.4.
Together they constitute the exterior Schwarzschild metric.

%% ============================================================
\section{Validation via the Bekenstein Bound}
%% ============================================================

The value $\alpha = 2\pi$ was fixed in Section~2.3 by matching to
Schwarzschild.
We now show it is independently confirmed by the Bekenstein bound.

\subsection{The Bekenstein Bound}

The Bekenstein bound limits the information content of a region of
radius $r$ containing energy $E = Mc^2$.
In natural units ($\hbar = c = 1$):
\begin{equation}
  K_{\max} = 2\pi Mr.
  \label{eq:bekenstein-bound}
\end{equation}
The factor $2\pi$ appears in Bekenstein's original derivation,
reflecting the relationship between energy, radius, and holographic
capacity.

\subsection{Information in Transit}

From Section~2.3, with $\alpha = 2\pi$, the information in transit
between the source and a sphere at radius $r$ (in natural units) is
\begin{equation}
  \Delta C = \frac{2\pi Mc^2}{\hbar}\cdot\frac{r}{c} = 2\pi Mr.
  \label{eq:transit-natural}
\end{equation}

\subsection{Exact Equality and Its Significance}

Comparing equations~\eqref{eq:bekenstein-bound}
and~\eqref{eq:transit-natural}: $\Delta C = K_{\max}$ exactly.
The information in transit at every radius saturates the Bekenstein
bound.
Three implications follow.

\medskip
\noindent\textbf{Independence.}
$\alpha = 2\pi$ was fixed by matching to Schwarzschild.
The Bekenstein bound provides a completely independent route to the
same value, including the exact factor of $2\pi$.
This is not a fitted parameter.

\medskip
\noindent\textbf{Saturation as a governing principle.}
Mass generates information at exactly the maximum rate compatible with
the geometric capacity of its surroundings.
This ``just-in-time'' character ensures consistency without fine-tuning.

\medskip
\noindent\textbf{Compton scale validation.}
The equality holds at all radii, including the Compton wavelength
$\lambda_C = \hbar/Mc$, providing a natural consistency check at the
quantum scale.

\subsection{Summary}

The rate $dC/d\tau = 2\pi Mc^2/\hbar$: reproduces the Schwarzschild
radial metric (Section~2.3); saturates the Bekenstein bound at every
radius (this section); yields the factor $2\pi$ from two independent
requirements.
These establish $\alpha = 2\pi$ as an inevitable consistency
requirement, not a free choice.

%% ============================================================
\section{Discussion and Conclusions}
%% ============================================================

\subsection{Summary of Results}

We have derived the radial component of the exterior Schwarzschild
metric from one hypothesis (H) and four assumptions (A1--A4), without
invoking Einstein's field equations.
The radial component $g_{rr}$ is exact.
The temporal component $g_{tt}$ rests on a well-motivated conjecture
identified in Section~2.4; an independent derivation appears in
Paper~II.
Matching fixes $\alpha = 2\pi$, independently confirmed by the
Bekenstein bound.

\subsection{Extensions in Companion Papers}

The following topics are developed in companion papers.

\medskip
\noindent\textbf{Einstein equations (Paper~II).}
The full Einstein field equations for vacuum and pressureless dust,
derived from Screen Theory premises.
$g_{tt}$ is recovered here as a theorem rather than a conjecture.

\medskip
\noindent\textbf{Cosmology (Paper~III).}
Global information balance, Machian origin of mass-energy, and the
cosmological constant reinterpreted in information-theoretic terms.

\medskip
\noindent\textbf{Horizons and black holes (Paper~IV).}
Behavior at and inside $r = r_S$, information export rates, and the
Page curve in information-theoretic terms.

\subsection{Conceptual Implications}
\label{sec:implications}

The derivation suggests a fundamental reinterpretation.
Mass is not a source of curvature in the sense of a cause producing an
effect; it \emph{is} a rate of information generation, and curvature
is the geometric shape that rate imposes on its surroundings.
This makes precise Wheeler's \cite{Wheeler1990} ``it from bit'' vision:
geometry (\emph{that}) from information generation (\emph{nat}).

The causal order is inverted relative to standard GR\@.
In GR, mass-energy is a source term in dynamical field equations and
curvature is the solution.
Here, irreversible information generation is primary and geometry is
its kinematic accommodation.
The Compton frequency $Mc^2/\hbar$, typically interpreted as a quantum
oscillation frequency, acquires new significance as the characteristic
rate of irreversible constraint production.

\medskip
\noindent\textbf{Newton's gravitational constant reinterpreted.}
From the derivation chain, $G$ can be written as
\begin{equation*}
  G = \frac{\ell_P^2 c^3}{\hbar},
\end{equation*}
where $4\ell_P^2$ converts area to information (A3) and $\hbar$
converts action to information.
$G$ is therefore not a fundamental coupling constant but an
\emph{exchange rate} between Planck area and quantum of action.
The factor $c^3$ adjusts units only.
$G$'s apparent smallness is not a mystery requiring explanation but
a consequence of the ratio between two independently motivated
fundamental quantities.

\medskip
\noindent\textbf{Proper radial length and holographic saturation.}
Once $g_{rr}$ is in hand, the ratio of coordinate displacement to
proper radial length follows immediately:
\begin{equation*}
  \frac{dr}{d\ell}
  = \frac{1}{\sqrt{g_{rr}}}
  = \sqrt{1 - \frac{r_S}{r}}
  = \sqrt{1 - \frac{S_{\rm Bekenstein}}{S_{\max}}}.
\end{equation*}
The ratio of coordinate to proper radial length is the square root
of the \emph{unsaturated} holographic fraction.
At $r = r_S$, saturation is complete: $dr/d\ell \to 0$, and an
infinite proper distance separates any exterior point from the horizon.
The horizon is the surface of complete holographic saturation, and the
divergence of $g_{rr}$ there is not a coordinate artifact but a direct
geometric expression of that saturation.

\subsection{An Open Problem and an Invitation}

The derivation of $g_{rr}$ is rigorous within the stated assumptions.
The argument for $g_{tt}$ rests on a step --- minimality excluding
unsourced free parameters --- that we present as a conjecture.
A rigorous proof of this step, a counterexample forcing modification
of the framework, or an alternative derivation closing the gap, would
each be a valuable contribution.
The companion paper (Paper~II) provides an independent route to $g_{tt}$
via the Einstein equations, which may serve as confirmation of the
conjecture pending a direct proof.

\subsection{Forward-Looking Conjectures}

\noindent\textbf{Parameter-free derivation.}
The Bekenstein bound, treated as an equality at the Compton scale,
combined with the generation rate hypothesis, may yield $\alpha = 2\pi$
with no matching to Schwarzschild required.

\medskip
\noindent\textbf{Shape independence.}
Bekenstein's bound assumes nothing about interior mass distribution.
The present argument may inherit this insensitivity; the spherically
symmetric case may be the simplest instance of something more general.

\medskip
\noindent\textbf{Screen overlap and Newton's law.}
Newtonian gravity and its relativistic extension may admit a unified
interpretation via holographic screen overlap between masses, with the
nonlinearity of Einstein's equations corresponding to the breakdown of
the small-overlap approximation.

\medskip
\noindent\textbf{Cosmological information accumulation.}
The accumulation of classical information on cosmic horizons --- the
irreversible deposit of settled outcomes beyond the reach of all future
observers --- may drive large-scale spacetime dynamics.
This connection, including the Machian origin of mass-energy and the
derivation of $E = Mc^2$ from Screen Theory premises, is developed in
Paper~III.

%% ============================================================
\section*{Acknowledgments}
%% ============================================================

The author thanks Joel Welling for illuminating early conversations
about the holographic principle, and Bob Carlitz for valuable feedback.

%% ============================================================
\section*{References}
%% ============================================================

\begin{enumerate}

  \item\label{Bekenstein1973}
        Bekenstein, J.\,D.\ (1973).
        Black holes and entropy.
        \textit{Physical Review D}, 7(8), 2333--2346.

  \item\label{Bekenstein1981}
        Bekenstein, J.\,D.\ (1981).
        Universal upper bound on the entropy-to-energy ratio for
        bounded systems.
        \textit{Physical Review D}, 23(2), 287--298.

  \item\label{HawkingPenrose1996}
        Hawking, S.\,W.\ and Penrose, R.\ (1996).
        \textit{The Nature of Space and Time}.
        Princeton University Press.

  \item\label{Kolmogorov1965}
        Kolmogorov, A.\,N.\ (1965).
        Three approaches to the quantitative definition of information.
        \textit{Problems of Information Transmission}, 1(1), 1--7.

  \item\label{Mach1883}
        Mach, E.\ (1883).
        \textit{Die Mechanik in ihrer Entwicklung historisch-kritisch
        dargestellt}.
        Brockhaus, Leipzig.
        [English translation: \textit{The Science of Mechanics},
        Open Court, 1893.]

  \item\label{PDG2024}
        Navas, S.\ et al.\ (Particle Data Group) (2024).
        \textit{Review of Particle Physics}.
        \textit{Physical Review D}, 110, 030001.

  \item\label{Penrose1994}
        Penrose, R.\ (1994).
        \textit{Shadows of the Mind}.
        Oxford University Press.

  \item\label{Sciama1953}
        Sciama, D.\,W.\ (1953).
        On the origin of inertia.
        \textit{Monthly Notices of the Royal Astronomical Society},
        113, 34--42.

  \item\label{Shannon1948}
        Shannon, C.\,E.\ (1948).
        A mathematical theory of communication.
        \textit{Bell System Technical Journal},
        27(3), 379--423 and 27(4), 623--656.

  \item\label{Srednicki1993}
        Srednicki, M.\ (1993).
        Entropy and area.
        \textit{Physical Review Letters}, 71(5), 666--669.

  \item\label{Susskind1995}
        Susskind, L.\ (1995).
        The world as a hologram.
        \textit{Journal of Mathematical Physics}, 36(11), 6377--6396.

  \item\label{tHooft1993}
        't~Hooft, G.\ (1993).
        Dimensional reduction in quantum gravity.
        arXiv:gr-qc/9310026.

  \item\label{Wald1984}
        Wald, R.\,M.\ (1984).
        \textit{General Relativity}.
        University of Chicago Press.

  \item\label{Wheeler1990}
        Wheeler, J.\,A.\ (1990).
        Information, physics, quantum: The search for links.
        In Zurek, W.\,H.\ (ed.),
        \textit{Complexity, Entropy, and the Physics of Information}.
        Addison-Wesley.

\end{enumerate}

%% ============================================================
\appendix
\section{Cosmological Motivation for Assumption A2}
\label{app:A2}
%% ============================================================

Assumption A2 --- that the vacuum near matter operates at holographic
saturation --- is the least self-evident of the four assumptions
underlying the derivation.
We offer here a quantitative motivation, deferring the exact
calculation for the concordance cosmology to Paper~III.

\subsection*{Intellectual context}

The question of where space comes from has a long history.
Leibniz argued against Newton that space has no existence independent
of the relations between objects; to posit an absolute space is to
posit something with no observable consequences and no causal role.
Mach \cite{Mach1883} extended this relational programme to inertia,
proposing that the resistance of matter to acceleration is not a
response to absolute space but to the gravitational influence of all
other matter in the universe.
General relativity absorbed Mach's relational spirit without fully
vindicating it: the metric is relational in form, but the manifold
remains a given substrate.

Sciama \cite{Sciama1953} gave Mach's proposal its most precise
quantitative form, showing that the sum of gravitational potential
contributions $\sum GM_i/r_i c^2$ from all matter in the observable
universe yields a result of order unity, suggesting that the inertia
of any local object is entirely due to its interaction with
distant matter.
The calculation below is structurally identical to Sciama's integral,
evaluated in information-theoretic language rather than gravitational
potential language.
The agreement is not a coincidence; both are expressions of the same
underlying physics.

\subsection*{The Olbers integral for holographic area}

We model the universe as Einstein--de~Sitter: spatially flat,
matter-dominated, $\Omega = 1$, $\Lambda = 0$, with critical density
$\rho_c = 3H_0^2/8\pi G$ and particle horizon $r_H = 2c/H_0$.
Rather than summing light flux from distant sources as in Olbers's
paradox, we sum holographic area contributions.

From Section~2.3, a mass $M$ at distance $r$ contributes information
in transit
\begin{equation}
  \Delta C = \frac{2\pi Mcr}{\hbar},
  \label{eq:app-transit}
\end{equation}
which by A3 requires holographic area
\begin{equation}
  \Delta A = 4\ell_P^2\,\Delta C = \frac{8\pi GMr}{c^2}.
  \label{eq:app-deltaA}
\end{equation}
Integrating over concentric shells at critical density:
\begin{equation}
  A_{\rm required}
  = \int_0^{r_H}\frac{8\pi Gr}{c^2}\,\rho_c\,4\pi r^2\,dr
  = \frac{128\pi^2 G\rho_c}{c^2}\cdot\frac{r_H^4}{4}
  = 48\pi\frac{c^2}{H_0^2}.
  \label{eq:app-required}
\end{equation}
The holographic capacity of the cosmological horizon is
\begin{equation}
  A_{\rm horizon} = 4\pi r_H^2 = 16\pi\frac{c^2}{H_0^2},
  \label{eq:app-horizon}
\end{equation}
giving the ratio
\begin{equation}
  \frac{A_{\rm required}}{A_{\rm horizon}} = 3.
  \label{eq:app-ratio}
\end{equation}
This is a factor-of-3 result with no free parameters.

By the Copernican principle this holds at every point: each location
receives information flux from its own observable universe, and that
flux saturates the local holographic screen to within a factor of~3
under the Einstein--de~Sitter approximation.
We conjecture that the exact calculation for the concordance cosmology
--- incorporating dark energy and the redshift evolution of matter
density --- closes this gap and yields exact saturation.
The factor of~3 overshoot in the EdS case is qualitatively consistent
with dark energy reducing the effective matter volume relative to
horizon area.

\subsection*{Physical interpretation}

This result provides independent quantitative motivation for A2:
the vacuum operates at holographic saturation not because we
postulate it, but because the collective information flux from all
matter in the observable universe fills the local holographic screen
to capacity at every point.

It also suggests a more complete resolution to the question Leibniz
and Mach were asking.
Space is not a pre-existing stage.
The local vacuum exists, and has the geometric structure it has,
because the universe has been generating information long enough and
densely enough that the past light cone of every point is saturated.
Spacetime is not ordained --- it is accumulated, built incrementally
from the information output of all matter within causal reach.

This reframes the apparent wastefulness of cosmic voids.
Every patch of apparently empty space encodes the complete history
of its past light cone, back to the beginning of time: a holographic
record of everything that has ever been causally accessible from
that location.
Matter is a light dusting of active information-generation centers
embedded in a universe that is overwhelmingly composed of this
accumulated record.

The Machian implication runs deeper still.
The same integral that gives equation~\eqref{eq:app-ratio} is
structurally identical to Sciama's \cite{Sciama1953} calculation of
inertial origin.
This suggests that mass-energy itself --- and ultimately $E = Mc^2$
--- has a Machian, relational origin: the energy content of matter
is built up from its interactions with the cosmological information
field rather than being intrinsic to matter in isolation.
Screen Theory may thus vindicate Mach's programme in a precise
quantitative sense: both space and inertia are accumulated,
relational quantities, not given substrates.
This connection is developed in Paper~III.

\subsection*{Note on placement}

This appendix presents the argument at the level of a toy calculation
sufficient to establish plausibility.
The full derivation --- with proper cosmological distance measure,
redshift corrections, and exact treatment of the concordance cosmology
--- is reserved for Paper~III, where it forms part of the central
argument.
This appendix may be read independently of the main text and may be
updated or replaced as the Paper~III calculation is completed.

\end{document}
